\answer{ex:19:1}{(а)~Один синапс дає $6.7\mV$ ($R=66.7$~МОм), тож до порога взагалі треба не менше $4$ (трьох вистачає лише асимптотично); $8$ за $10\ms$; $14$ за $5\ms$. (б)~$0.1\ms\ll\tau_m$, поправка від витоку $\sim0.25\%$.}
\answer{ex:19:2}{(а)~$1.75\cdot10^{10}$ (чверть змішувачів відповідає на все підряд), $4.4\cdot10^9$ і $0.06$ (за $k=40$ не відповідає майже ніколи). (б)~У $10^3$ разів: $2\cdot10^5$ проти $200$.}
\answer{ex:19:3}{(а)~$C(2n,n)=2^{2n-1}$ із симетрії біноміальних коефіцієнтів. (б)~$0.93$ і $0.09$; ширина переходу $\sim\sqrt{2n}$ за $P$, відносна $\sim1/\sqrt n$.}
\answer{ex:19:4}{З одного дендрита $V_{\text{тіло}}<\kappa E_E<\theta$ за будь-якої кількості входів; за $g_0=0.5g_L$ потрібно $n\ge3$ на кожному дендриті.}
\answer{ex:19:5}{$I\ge C\sigma_V/\sigma_t=15$~нА, тобто $150$ синхронних синапсів --- нереалістично; маленькому нейронові вистачило б $1$~нА, а із синапсом у $4$~нА тремтіння $5$~мкс.}
\answer{ex:19:6}{Максимум на ближньому кінці: $0.03\,\tau_m$ і $0.97$ ($L=0.2$); $0.32\,\tau_m$ і $0.66$ ($L=1$); $0.79\,\tau_m$ і $0.33$ ($L=2$). Оцінка~\eqref{eq:dendriteload} за повною ємністю правильна лише за $L\ll1$.}
\answer{ex:19:7}{Відкрита задача. За фіксованого $m$ час відгуку зростає лінійно із запам'ятовуваним $P$; за фіксованої частки $m=\varphi n$ він перестає залежати від $n$, і повільність великого нейрона --- наслідок підсумовування багатьох входів, а не розміру як такого.}
